\chapter{Introduction}\label{r02:ch:introduction}

waLBerla~\cite{walberla} is a massively parallel C++ framework for lattice Boltzmann
and other grid-based simulations. It runs on laptops and on the largest
supercomputers, on x86 and ARM CPUs and on NVIDIA and AMD GPUs. Its fastest kernels
are not written by hand: they are generated from a symbolic description by
\ps{}~\cite{pystencils}, during the build.

This report explains how waLBerla was made to run with explicit RISC-V Vector (RVV)
code on the SpacemiT~K1, the eight-core RISC-V chip of the Banana~Pi BPI-F3 board.
It is written for a developer who knows C++, Python and SIMD programming, but has
never looked inside waLBerla. The first part therefore introduces waLBerla itself: the
problem it solves, how its sources and its build system are organized, and how its
code generation works. The second part identifies what was missing for RVV, and
describes what was added, why it was done that way, how to build and run it, and what
it achieves.

\section{The short version}

\begin{keyidea}[The one thing to understand first]
  waLBerla never emits SIMD code itself. Its performance-critical loops (the
  \term{sweeps}) are C++ files that the code generator \ps{} writes at build time.
  So ``RVV support in waLBerla'' consists of three parts: a RISC-V vector
  \emph{target} in \ps{}; waLBerla \emph{build flags} that tell the C++ compiler
  about the vector unit; and a generator script that \emph{asks} for vectorized
  code. Each part was missing.
\end{keyidea}

Before this work, waLBerla on the K1 behaved as follows (\cref{r02:ch:gap}):
\begin{itemize}
  \item With the default settings it built, but every kernel was scalar: the compiler
        was not told about the vector extension, and sweepgen generated plain C++.
  \item With \cmakeopt{WALBERLA_OPTIMIZE_FOR_LOCALHOST=ON}, the switch that selects
        the host's best instruction set on x86 and ARM, it did not build at all. The
        option adds \code{-march=native}, which both GCC~14 and Clang~18 reject on
        RISC-V. And had it built, \ps{} would still have generated plain C++: it knew no
        RISC-V vector target, and its host detection only warned ``unknown platform
        riscv64''.
  \item Even on x86, waLBerla's new code generation layer (sweepgen) never turns
        vectorization on. A generator script has to request it.
  \item Around these gaps were smaller obstacles: Python package versions that did not
        match the pinned \ps{}, a C++20 feature in the halo exchange that Clang~18 does
        not implement, and no benchmark that could compare the variants.
\end{itemize}

What was added (\cref{r02:ch:design,r02:ch:pystencils,r02:ch:wlb-changes}):
\begin{itemize}
  \item \textbf{In \ps{}} (fork commit \texttt{9a989e16}, about 1270 lines including
        tests): the target \code{Target.RISCV_RVV}, host detection from
        \file{/proc/cpuinfo}, a platform that maps the vector IR to RVV~1.0 C intrinsics
        for a fixed register width of 256 bits, a small runtime header, and the matching
        JIT compiler flags. The design copies the contract of the existing AVX backend
        on purpose.
  \item \textbf{In waLBerla} (fork commits \texttt{dbb1b400} and \texttt{b78f2d87}): an
        explicit \code{-march=rv64gcv_zvl256b} instead of \code{-march=native} on
        riscv64 hosts, for GCC and Clang (about 20 lines of CMake); a fix that lets Clang~18
        compile the v8 halo exchange; and a new benchmark application,
        \file{apps/benchmarks/SweepgenLbmCPU}, that builds the same LBM sweep as
        scalar, auto-vectorized and explicit-RVV code. \emph{No new CMake option and no
        change to sweepgen} were needed: the existing path
        \cmakeopt{WALBERLA_OPTIMIZE_FOR_LOCALHOST} $\rightarrow$ \code{Target.CurrentCPU}
        now resolves to \code{RISCV_RVV} on the K1, exactly as it resolves to AVX on
        an x86 workstation.
  \item \textbf{In this workspace}: board scripts that build and run the benchmark
        matrix unattended, and the measurements in \cref{r02:ch:results}.
\end{itemize}

The result: on one core of the K1, the explicit RVV sweep of a D3Q19 lattice
Boltzmann simulation is 1.16--1.23$\times$ faster than the best auto-vectorized build in
double precision and 1.50--1.70$\times$ faster in single precision. The checksums of all
double-precision variants agree to twelve digits. With all eight cores the gain shrinks
to 1.05--1.26$\times$, because the sweep then runs into the memory bandwidth of the
chip.

\section{How to read this report}

\begin{description}[style=nextline]
  \item[\Cref{r02:ch:walberla}: waLBerla in a nutshell] The lattice Boltzmann
    method in one page, the block-structured domain, fields, sweeps, halo exchange and
    time loop, and the difference between the classic core and the new \emph{v8} core.
  \item[\Cref{r02:ch:structure}: Project structure and build system] The source tree,
    the CMake options that matter here, and where compiler flags come from.
  \item[\Cref{r02:ch:codegen}: Code generation] How a Python generator script becomes
    a C++ sweep class during the build, and where the code generator target is chosen.
  \item[\Cref{r02:ch:rvv-primer}: RVV and the \ps{} backend] The RVV concepts needed
    for the rest of the report, and what the generated code looks like.
  \item[\Cref{r02:ch:gap,r02:ch:design}] What was missing, layer by layer, and the
    design decisions (including a first, larger attempt that was later discarded).
  \item[\Cref{r02:ch:pystencils,r02:ch:wlb-changes}] The changes in both forks, file by
    file.
  \item[\Cref{r02:ch:build-run,r02:ch:results,r02:ch:limitations}] How to build and run
    it on the board, validation and benchmark results, and known limitations.
\end{description}
Readers who already know waLBerla can start at \cref{r02:ch:gap}. Report~1 of this
collection is a manual of \ps{} itself; its \cref{r01:ch:rvv-backend} documents the
RVV backend in more depth than \cref{r02:ch:pystencils} of this report does.

\section{Sources and versions}\label{r02:sec:versions}

Everything described here lives in this workspace (\file{my-e4-ws}):
\begin{center}\small
\begin{tabularx}{\linewidth}{@{}l>{\raggedright\arraybackslash}X@{}}
  \toprule
  \textbf{Component} & \textbf{Version used}\\
  \midrule
  waLBerla & \file{submodules/walberla}, upstream \texttt{v7.2-88-g0c8ed8c9}
    (2026-05) plus three fork commits on branch \texttt{rvv-support}:
    \texttt{e9712d59} (first attempt, superseded), \texttt{dbb1b400} (build flags),
    \texttt{b78f2d87} (benchmark app, Clang fix)\\
  \ps{} & \file{submodules/pystencils}, upstream \texttt{20d7dcf0} (2.0b1) plus
    \texttt{31758e93} (first attempt, superseded) and \texttt{9a989e16} (the VLS-256
    backend) on branch \texttt{rvv-support}\\
  pystencils-sfg & 0.1.1 (PyPI), the last release that works with the pinned \ps{}\\
  lbmpy & 2.1.1 (PyPI, installed without dependencies)\\
  Board & Banana~Pi BPI-F3: SpacemiT~K1, 8 SpacemiT~X60 cores in two clusters of 4,
    RVV~1.0 with VLEN = 256 bits, 4\,GB RAM; Bianbu Linux (Ubuntu 24.04 based)\\
  Toolchains & GCC 14.2.0 and Clang 18.1.8 (both Bianbu builds), CMake 3.28.3,
    Open MPI 4.1.6, Python 3.12\\
  \bottomrule
\end{tabularx}
\end{center}
The design notes and the full measurement log are in
\file{docs/notes2/rvv-vls256-only-plan-opus5.5.md} and
\file{docs/notes2/rvv-vls256-only-progress-opus5.5.md}. An older, broader analysis of
waLBerla's code generation paths is \file{docs/notes/walberla-codegen-flow.md}. Paths in
this report are relative to the waLBerla or \ps{} submodule root unless they start with
\file{submodules/}, \file{scripts/} or \file{docs/}, which are relative to the workspace
root.
